
\documentclass{article}
\usepackage{colortbl}
\usepackage{makecell}
\usepackage{multirow}
\usepackage{supertabular}

\begin{document}

\newcounter{utterance}

\centering \large Interaction Transcript for game `eqbench', experiment `eq\_bench', episode 43 with reflexion{-}llama{-}3.1{-}8b{-}instruct.
\vspace{24pt}

{ \footnotesize  \setcounter{utterance}{1}
\setlength{\tabcolsep}{0pt}
\begin{supertabular}{c@{$\;$}|p{.15\linewidth}@{}p{.15\linewidth}p{.15\linewidth}p{.15\linewidth}p{.15\linewidth}p{.15\linewidth}}
    \# & \multicolumn{2}{c}{Player} && \multicolumn{2}{c}{Game Master} \\
    \hline

    \theutterance \stepcounter{utterance}  
    & & & \multicolumn{4}{p{0.6\linewidth}}{
        \cellcolor[rgb]{0.9,0.9,0.9}{
            \makecell[{{p{\linewidth}}}]{
                \texttt{\tiny{[P1$\langle$GM]}}
                \texttt{You have attempted this same task before and did not succeed. These are the plans you wrote after those attempts.} \\
\texttt{Trial \#1: New plan: } \\
\\ 
\texttt{Upon re{-}reading the dialogue, I realized that I focused too much on the surface{-}level conversation and didn't delve deeper into the emotional dynamics between James and Rachel. I should have paid more attention to the underlying tensions and power struggles in their relationship. } \\
\\ 
\texttt{In particular, I should have noticed that Rachel is trying to push James out of his comfort zone and make him confront his feelings, while James is struggling to express himself and is feeling intimidated by Rachel's emotional intensity. } \\
\\ 
\texttt{To improve my analysis, I will focus on identifying the emotional triggers and motivations behind each character's actions and words. I will also pay closer attention to the power dynamics and the ways in which the characters are influencing each other's emotions. } \\
\\ 
\texttt{For example, I might ask myself questions like: What is driving Rachel's determination to rehearse the kiss scene? Is it a genuine desire to improve the play, or is she trying to push James out of his comfort zone? How does James's fear of intimacy and his tendency to retreat into his shell affect his interactions with Rachel? } \\
\\ 
\texttt{By taking a more nuanced and empathetic approach, I hope to improve my understanding of the emotional dynamics at play and provide a more accurate analysis of Rachel's emotional state at the end of the dialogue.} \\
\\ 
\\ 
\\ 
\texttt{Your task is to predict the likely emotional responses of a character in this dialogue:} \\
\\ 
\texttt{James: I can't believe we're doing this again, Rachel. } \\
\texttt{Rachel: What, the play? Or you mean us, playing lovers?} \\
\texttt{James: Both, I guess. } \\
\texttt{Rachel: Well, it's not like we have a choice. We were cast for these roles. } \\
\texttt{James: Yeah, but it's... complicated.} \\
\texttt{Rachel: Complicated? Are you still stuck on that scene from three years ago?} \\
\texttt{James: I wouldn't say stuck. More like... haunted.} \\
\texttt{Rachel: Haunted? That's a bit dramatic, even for you, James.} \\
\texttt{James: Maybe, but it's the truth.} \\
\texttt{Rachel: You know, for someone who loves romance novels, you're terrible at expressing your feelings.} \\
\texttt{James: And you're too good at it. It's intimidating.} \\
\texttt{Rachel: Intimidating? Me? You're the one who's always retreating into your shell.} \\
\texttt{James: It's not like I do it on purpose, Rachel. It's just... hard.} \\
\texttt{Rachel: What's hard, James? Acting or being around me?} \\
\texttt{James: Both.} \\
\texttt{Rachel: Well, we're going to have to figure it out, aren't we? For the play.} \\
\texttt{James: Yes, for the play.} \\
\texttt{Rachel: So, let's rehearse the scene. The kiss scene.} \\
\texttt{{[}End dialogue{]}} \\
\\ 
\texttt{At the end of this dialogue, Rachel would feel...} \\
\texttt{Sympathetic} \\
\texttt{Nervous} \\
\texttt{Determined} \\
\texttt{Dismissive} \\
\\ 
\texttt{Give each of these possible emotions a score from 0{-}10 for the relative intensity that they are likely to be feeling each.} \\
\\ 
\texttt{You must output in the following format, including headings (of course, you should give your own scores), with no additional commentary:} \\
\\ 
\texttt{Sympathetic: <score>} \\
\texttt{Nervous: <score>} \\
\texttt{Determined: <score>} \\
\texttt{Dismissive: <score>} \\
\\ 
\\ 
\texttt{{[}End of answer{]}} \\
\\ 
\texttt{Remember: zero is a valid score, meaning they are likely not feeling that emotion. You must score at least one emotion > 0.} \\
\\ 
\texttt{Your answer:} \\
            }
        }
    }
     \\ \\

    \theutterance \stepcounter{utterance}  
    & \multicolumn{4}{p{0.6\linewidth}}{
        \cellcolor[rgb]{0.561,0.737,0.561}{
            \makecell[{{p{\linewidth}}}]{
                \texttt{\tiny{[P1$\rangle$GM]}}
                \texttt{Sympathetic: 8} \\
\texttt{Nervous: 6} \\
\texttt{Determined: 9} \\
\texttt{Dismissive: 0} \\
            }
        }
    }
    & & \\ \\

    \theutterance \stepcounter{utterance}  
    & & & \multicolumn{2}{p{0.3\linewidth}}{
        \cellcolor[rgb]{0.9,0.9,0.9}{
            \makecell[{{p{\linewidth}}}]{
                \texttt{\tiny{[GM$|$GM]}}
                \texttt{Sympathetic: 8} \\
\texttt{Nervous: 6} \\
\texttt{Determined: 9} \\
\texttt{Dismissive: 0} \\
            }
        }
    }
    & & \\ \\

    \theutterance \stepcounter{utterance}  
    & & & \multicolumn{2}{p{0.3\linewidth}}{
        \cellcolor[rgb]{0.9,0.9,0.9}{
            \makecell[{{p{\linewidth}}}]{
                \texttt{\tiny{[GM$|$GM]}}
                \texttt{Sympathetic: 4} \\
\texttt{Nervous: 0} \\
\texttt{Determined: 8} \\
\texttt{Dismissive: 2} \\
            }
        }
    }
    & & \\ \\

    \theutterance \stepcounter{utterance}  
    & & & \multicolumn{2}{p{0.3\linewidth}}{
        \cellcolor[rgb]{0.9,0.9,0.9}{
            \makecell[{{p{\linewidth}}}]{
                \texttt{\tiny{[GM$|$GM]}}
                \texttt{game\_result = LOSE} \\
            }
        }
    }
    & & \\ \\

    \theutterance \stepcounter{utterance}  
    & & \multicolumn{4}{p{0.6\linewidth}}{
        \cellcolor[rgb]{0.95,0.95,0.95}{
            \makecell[{{p{\linewidth}}}]{
                \texttt{\tiny{[P1$\leftrightarrow$P]}}
                \texttt{Post{-}episode Reflexion (not sent to GM):} \\
\texttt{New plan: } \\
\\ 
\texttt{Upon re{-}reading the dialogue, I realized that I focused too much on the surface{-}level conversation and didn't delve deeper into the emotional dynamics between James and Rachel. I should have paid more attention to the underlying tensions and power struggles in their relationship. } \\
\\ 
\texttt{In particular, I should have noticed that Rachel is trying to push James out of his comfort zone and make him confront his feelings, while James is struggling to express himself and is feeling intimidated by Rachel's emotional intensity. } \\
\\ 
\texttt{To improve my analysis, I will focus on identifying the emotional triggers and motivations behind each character's actions and words. I will also pay closer attention to the power dynamics and the ways in which the characters are influencing each other's emotions. } \\
\\ 
\texttt{For example, I might ask myself questions like: What is driving Rachel's determination to rehearse the kiss scene? Is it a genuine desire to improve the play, or is she trying to push James out of his comfort zone? How does James's fear of intimacy and his tendency to retreat into his shell affect his interactions with Rachel? } \\
\\ 
\texttt{I will also consider the following key points:} \\
\\ 
\texttt{* Rachel's comment about James being "terrible at expressing his feelings" suggests that she is trying to provoke him into opening up.} \\
\texttt{* James's admission that he is "haunted" by a past scene implies that he is struggling with emotional baggage.} \\
\texttt{* Rachel's response to James's fear of intimacy, saying that he is "too good at it" and that it's "intimidating," suggests that she is trying to push him out of his comfort zone.} \\
\\ 
\texttt{By taking a more nuanced and empathetic approach, I hope to improve my understanding of the emotional dynamics at play and provide a more accurate analysis of Rachel's emotional state at the end of the dialogue.} \\
\\ 
\texttt{Sympathetic: 8} \\
\texttt{Nervous: 6} \\
\texttt{Determined: 9} \\
\texttt{Dismissive: 0} \\
            }
        }
    }
    & \\ \\

\end{supertabular}
}

\end{document}
